\documentclass{beamer}

\mode<presentation>
{
  \usetheme{Frankfurt}      
  \usecolortheme{beaver}    
  \usefonttheme{default}
  \setbeamertemplate{navigation symbols}{}
  \setbeamertemplate{caption}[numbered]
} 

\usepackage[english]{babel}
\usepackage[utf8]{inputenc}
\usepackage{graphicx}

% source displacement, gaussian distribution, plots with different parameter combinations

\logo{\includegraphics[height=1cm]{logo_unipv.png}}

\title[1D Tracker Project]{Digital Data Acquisition Techniques: \\ Reconstruction of a 1D tracker}
\author{Pietro Schiavone}
\institute[UNIPV]{University of Pavia}
\date{\today}

\begin{document}

\begin{frame}
  \titlepage
\end{frame}

\begin{frame}{Overview}
  \tableofcontents
\end{frame}

% ==========================================
\section{Goals}
% ==========================================

\begin{frame}{Goals \& Working Assumptions}
\begin{itemize}
    \item \textbf{Goals:} Reconstruction of the trajectories of charged particles crossing N layers of a one-dimensional tracker.
    \item \textbf{Setup:}
    \begin{itemize}
        \item $N$ equally spaced layers equipped with boolean sensors distributed continuously along each layer.
        \item Point-like source, with emission angles generated from specific probability distributions (e.g. Uniform and Gaussian).
        \item Each layer has an acceptance of $[0, 1]$.
    \end{itemize}
\end{itemize}
\end{frame}

\begin{frame}{Experimental Set-up}
  \begin{center}
    \includegraphics[height=0.65\textheight]{apparato_sperimentale.pdf}
    
    \vspace{0.3cm}
    \scriptsize \textcolor{gray}{Simulation of 100 charged particles crossing 3 detector layers, each made of 400 pixels.}
  \end{center}
\end{frame}
% ==========================================
\section{Design, Demonstration and Validation}
% ==========================================

\begin{frame}{Design}
\textbf{Separation into Work Modules}:
    \begin{itemize}
        \item \textbf{Simulation:} Generates the tracks and produces the detector hits, which are then saved to disk using a specific Data Format.
        \item \textbf{Track reconstruction:} Given the hit detector pixels, estimates the tracks through a $\chi^2$ fit, using both the GPU and the CPU.
        \item \textbf{Analysis:} Reconstruction of the source position and filtering of the found tracks.
        \item \textbf{Evaluation:} Assessment of CPU/GPU performance and of the accuracy of the track-finding algorithm.
    \end{itemize}
\end{frame}

\begin{frame}{Object Oriented Architecture}
     \begin{itemize}
        \item \texttt{Source}: Sets the particle emission parameters and computes the trajectories.
        \item \texttt{Detector}: Converts hit positions into the corresponding pixels and vice versa.
        \item \texttt{Storage}: Handles the I/O operations and the encoding of the Event Format.
        \item \texttt{Config}: Main wrapper that manages the generation of event loops.
        \item \texttt{Visualization}: Handles the generation of the plots.
        \item \texttt{Truth}: Analysis for accuracy estimates.
    \end{itemize}
\end{frame}

\begin{frame}{Event Format}
\begin{itemize}
    \item Structured as an array of 32-bit words in binary format:
    \begin{itemize}
        \item \textbf{Event Header:} Contains the full Event Header \texttt{0xCACCAAAA}.
        \item \textbf{Metadata:} Following the event header come the metadata, including the number of layers, the length of each tracker layer, the pixel length, the inter-layer distance, the source-to-layer distance, the number of true tracks and the version.
        \item \textbf{Layer Header:} Each layer has its own checkword (\texttt{0xAAA0000n}) where n is the layer identifier, and each bit of the 32-bit words represents a detector pixel.
    \end{itemize}
\end{itemize}
\end{frame}

\begin{frame}{MC Truth}
At the end of the event format we find:
\begin{itemize}
    \item \textbf{MC checkword:} The checkword of the Monte Carlo data is \texttt{0xDEFAC0CA}.
    \item The particle emission angles from the source, so that the true tracks can be reconstructed.
\end{itemize}
\end{frame}

\begin{frame}{Example of the on-disk Data Format}
\begin{center}
    \includegraphics[height=0.8\textheight]{data_format.png}
\end{center}
\end{frame}
\begin{frame}{CPU Track Reconstruction: $\chi^2$ minimisation}
\begin{itemize}
    \item Iterating over the first layer, the algorithm computes the theoretical projections compatible with the detector acceptance.
    \item It then uses a binary search to select, on the last layer, only the contiguous portion of valid hits.
    \item The distance between the intermediate hits and the theoretical trajectory generated by the seed is scanned.
    \item A least-squares fit is performed on the final coordinates. The track is validated and stored only if the $\chi^2$, weighted on the pixel resolution, falls within the threshold.
\end{itemize}
 
\end{frame}
\begin{frame}{GPU Track Reconstruction: $\chi^2$ minimisation}
\begin{itemize}
    \item 2D Thread Mapping: the kernel assigns each single thread to an exclusive combination of one first-layer hit and one last-layer hit (using $16 \times 16$ thread blocks).
    \item Geometric Early-Exit.
    \item The GPU performs a sequential linear search over the intermediate layers to avoid thread divergence.
    \item Storage of the tracks in shared memory with AtomicAdd.
    \item The data are written to VRAM buffers allocated once, persistently, with no allocation/deallocation per event.
\end{itemize}
\end{frame}

\begin{frame}{Example:}
    \begin{center}
        \includegraphics[height=0.75\textheight]{apparato_combinatorio.pdf}
            
        \vspace{0.3cm}
        \scriptsize \textcolor{gray}{Raw tracks found (they are the same for both CPU and GPU) without any cut on the track q. N true tracks = 30; pixel length = 0.0025 m; Noise 1\%;}
    \end{center}
\end{frame}

\begin{frame}{Track reconstruction: ghost-track removal}
The track\_finder() function can reconstruct the same physical particle multiple times, with slightly different parameters, when a noise hit close to the true one still passes the $\chi^2$ cut.
\begin{itemize}
    \item RimuoviCloni() identifies duplicates with two combined criteria: (m,q) parameters compatible within 3$\sigma$ (the same theoretical sigmas propagated from the geometry) and at least one hit shared between the two reconstructions.
    \item The identified clones are merged into a single track: parameters = cluster average, $\chi^2$ = the best among the members.
\end{itemize}

\end{frame}

\begin{frame}{Example:}
      \begin{center}
        \includegraphics[height=0.75\textheight]{rimozione_cloni.pdf}
            
        \vspace{0.3cm}
        \scriptsize \textcolor{gray}{Removal of clone tracks from the raw tracks, with selection of tracks having q $\in [0.4, 0.6]$. N true tracks = 30; pixel length = 0.0025 m; Noise 1\%;}
    \end{center}  
\end{frame}

\begin{frame}{Example:}
    \begin{columns}[T] % The [T] option top-aligns the columns
    
        % --- FIRST COLUMN ---
        \begin{column}{0.5\textwidth}
            \begin{center}
                % Width adapted to fit within half a slide
                \includegraphics[width=\textwidth, keepaspectratio]{distrotracce_noq.pdf}
                
                \vspace{0.3cm}
                \scriptsize \textcolor{gray}{N true tracks = 30; pixel length = 0.0025 m; Noise 1\%;}
            \end{center} 
        \end{column}
        
        % --- SECOND COLUMN ---
        \begin{column}{0.5\textwidth}
            \begin{center}
                % Width adapted to fit within half a slide
                \includegraphics[width=\textwidth, keepaspectratio]{distrotracce.pdf}
                
                \vspace{0.3cm}
                \scriptsize \textcolor{gray}{N true tracks = 30; pixel length = 0.0025 m; Noise 1\%;}
            \end{center} 
        \end{column}
        
    \end{columns}
\end{frame}

\begin{frame}{Track reconstruction: source position}
Why not a classic gaussian fit?
    \begin{itemize}
        \item The hit positions are quantised and produce a peak made of discrete lines, not a continuous gaussian.
        \item A least-squares fit on this structure is unstable: it can collapse onto a single tooth, with $\sigma \rightarrow 0$.
    \end{itemize}
\end{frame}

\begin{frame}{Solution}
A deterministic two-stage estimate, with no non-linear minimisation.
    \begin{itemize}
        \item Stage 1: adaptive bin aggregation until the FWHM (full width at half maximum) reaches a physically credible value with respect to the theoretical sigma expected from the detector geometry $\implies$ obtain $\sigma$.
        \item Stage 2: iterative sigma-clipping with a fixed-width window, which only recentres the mean ($\mu$).
    \end{itemize}
\end{frame}

\begin{frame}{Source reconstruction}
  \begin{center}
    \includegraphics[height=0.5\textheight]{fit_gaussiano.pdf}
    
    \vspace{0.3cm}
    \scriptsize \textcolor{gray}{100000 calibration events were used for each N-tracks value, and efficiency and purity were then evaluated.}
\end{center}
\end{frame}

\begin{frame}{Fun fact}
If clones are not removed and the graphical binning is increased, one notices...
  \begin{center}
    \includegraphics[height=0.65\textheight]{pettine.pdf}
    
    \vspace{0.3cm}
    \scriptsize \textcolor{gray}{}
  \end{center}
\end{frame}

\begin{frame}{Why the comb structure?}
\begin{itemize}
    \item The detector hits are not continuous positions: each layer discretises them to the centre of the hit pixel.
    \item The reconstructed intercept q is a fixed-coefficient linear combination of the positions on the three layers (least-squares fit): being linear, it maps a lattice of discrete inputs into a discrete lattice of outputs.
\end{itemize}
\end{frame}
 
\begin{frame}{Example:}
    \begin{center}
        \includegraphics[height=0.75\textheight]{calibrazione.pdf}
            
        \vspace{0.3cm}
        \scriptsize \textcolor{gray}{Tracks remaining after calibration. N true tracks = 30; pixel length = 0.0025 m; Noise 1\%;}
    \end{center} 
\end{frame}

\begin{frame}{Noise}
\begin{itemize}
    \item Noise is added pixel by pixel: each still-off pixel has a fixed probability of turning on anyway, with an independent Bernoulli trial.
    \item Noise never switches off a true hit and does not alter pixels already lit by a particle.
    \item It is applied separately for each layer: no correlation between the noise of different layers.
\end{itemize}
\end{frame}

\begin{frame}{Example:}
    \begin{columns}[T] % The [T] option top-aligns the columns
    
        % --- FIRST COLUMN ---
        \begin{column}{0.5\textwidth}
            \begin{center}
                % Width adapted to fit within half a slide
                \includegraphics[width=\textwidth, keepaspectratio]{rumore_nofiltri.pdf}
                
                \vspace{0.3cm}
                \scriptsize \textcolor{gray}{N true tracks = 0; pixel length = 0.00025 m; Noise 3\%;}
            \end{center} 
        \end{column}
        
        % --- SECOND COLUMN ---
        \begin{column}{0.5\textwidth}
            \begin{center}
                % Width adapted to fit within half a slide
                \includegraphics[width=\textwidth, keepaspectratio]{rumore_rimozione_cloni.pdf}
                
                \vspace{0.3cm}
                \scriptsize \textcolor{gray}{N true tracks = 0; pixel length = 0.00025 m; Noise 3\%;}
            \end{center} 
        \end{column}
        
    \end{columns}
\end{frame}

\begin{frame}{Example:}
    \begin{center}
        \includegraphics[height=0.75\textheight]{rumore_fit.pdf}
            
        \vspace{0.3cm}
        \scriptsize \textcolor{gray}{100k runs with N true tracks = 0; pixel length = 0.00025 m; Noise 3\%;}
    \end{center}
\end{frame}

\begin{frame}{Efficiency and purity vs noise}
  \begin{center}
    \includegraphics[height=0.5\textheight]{vs_rumore.pdf}
    
    \vspace{0.3cm}
    \scriptsize \textcolor{gray}{100000 calibration events were used for each sampled noise percentage, and efficiency and purity were then evaluated.}
\end{center}
\end{frame}

\begin{frame}{Efficiency and purity vs N tracks}
  \begin{center}
    \includegraphics[height=0.5\textheight]{vs_tracce_generate.pdf}
    
    \vspace{0.3cm}
    \scriptsize \textcolor{gray}{100000 calibration events were used for each N-tracks value, and efficiency and purity were then evaluated.}
\end{center}
\end{frame}

\begin{frame}{WARNING}
in the efficiency and purity computation I did not account for all the tracks that exited the detector!
    \begin{center}
        \includegraphics[height=0.5\textheight]{tracce_perse.pdf}
    \end{center}

\end{frame}

\begin{frame}{WARNING}
    \begin{center}
        \includegraphics[height=0.5\textheight]{tracce_perse_vs_N.pdf}
    \end{center}
\end{frame}

\begin{frame}{Execution times}
  \begin{center}
    \includegraphics[height=0.65\textheight]{times.pdf}
    
    \vspace{0.3cm}
    \scriptsize \textcolor{gray}{Both curves measure the raw track search only. The GPU time uses the persistent-memory engine (GpuSingleEventFinder): there remains an irreducible fixed cost per event (kernel launch + PCIe transfers) O(1) $\mu s$.}
  \end{center}
\end{frame}

% ==========================================
\section{Next steps}
% ==========================================

\begin{frame}{Conclusions and Next Steps}
\begin{itemize}
    \item \textbf{Achievements:}
    \begin{itemize}
        \item  Implementation of a 1D tracker.
        \item  Efficient and validated event format.
        \item Track reconstruction with $\chi^2$ minimisation and reconstruction of the source position.
        \item Optimisation of GPU performance.
    \end{itemize}
    \item \textbf{Next Steps:}
    \begin{itemize}
        \item 2D/3D tracker.
        \item Layer misalignment.
        \item Introduction of dead space between pixels (non-continuous pixels).
        \item Track recognition with the Hough transform on FPGA, GPU-FPGA timing comparison, and comparison between the combinatorial method and the Hough transform.
    \end{itemize}
\end{itemize}
\end{frame}

\section{Backup}
\begin{frame}{Tracker with multiple sources}
	\begin{center}
        	\includegraphics[height=0.75\textheight]{multi_sorgente_solo_q_range.pdf}
        	\vspace{0.3cm}
        	\scriptsize \textcolor{gray}{N true particles = 5 / source; Pixel length = 0.0025 m; Noise = 3\%}
    	\end{center}
\end{frame}
\begin{frame}{Tracker with multiple sources}
	\begin{center}
        	\includegraphics[height=0.75\textheight]{multi_sorgente_cloni.pdf}
        	\vspace{0.3cm}
        	\scriptsize \textcolor{gray}{N true particles = 5 / source; Pixel length = 0.0025 m; Noise = 3\%}
    	\end{center}
\end{frame}
\begin{frame}{Tracker with multiple sources}
	\begin{center}
        	\includegraphics[height=0.75\textheight]{multi_sorgente_calibrato.pdf}        	        	
        	\vspace{0.3cm}
        	\scriptsize \textcolor{gray}{N true particles = 5 / source; Pixel length = 0.0025 m; Noise = 3\%}
    	\end{center}
\end{frame}
\begin{frame}{Tracker with multiple sources}
	\begin{center}
        	\includegraphics[height=0.75\textheight]{multi_sorgnete_distro_no_q.pdf}
        	\vspace{0.3cm}
        	\scriptsize \textcolor{gray}{N true particles = 5 / source; Pixel length = 0.0025 m; Noise = 3\%}
    	\end{center}
\end{frame}
\begin{frame}{Tracker with multiple sources}
	\begin{center}
        	\includegraphics[height=0.75\textheight]{multi_sorgente_distro_calibrato.pdf}
        	\vspace{0.3cm}
        	\scriptsize \textcolor{gray}{N true particles = 5 / source; Pixel length = 0.0025 m; Noise = 3\%}
    	\end{center}
\end{frame}
\begin{frame}{Tracker with multiple sources}
	\begin{center}
        	\includegraphics[height=0.75\textheight]{multi_sorgente_calibrazione_no_fit.pdf}
        	\vspace{0.3cm}
        	\scriptsize \textcolor{gray}{Calibration with 100k events. N true particles = 5 / source; Pixel length = 0.0025 m; Noise = 3\%}
    	\end{center}
\end{frame}
\begin{frame}{Tracker with multiple sources}
	\begin{center}
        	\includegraphics[height=0.75\textheight]{multi_sorgente_fit.pdf}
        	\vspace{0.3cm}
        	\scriptsize \textcolor{gray}{Calibration with 100k events. N true particles = 5 / source; Pixel length = 0.0025 m; Noise = 3\%}
    	\end{center}
\end{frame}
\begin{frame}{Displaced source}
	\begin{center}
        	\includegraphics[height=0.75\textheight]{spostata_no_filtri.pdf}
        	\vspace{0.3cm}
        	\scriptsize \textcolor{gray}{N true particles = 10; Pixel length = 0.00025 m; Noise = 3\%}
    	\end{center}    
\end{frame}
\begin{frame}{Displaced source}
	\begin{center}
        	\includegraphics[height=0.75\textheight]{spostata_no_cloni.pdf}
        	\vspace{0.3cm}
        	\scriptsize \textcolor{gray}{N true particles = 10; Pixel length = 0.00025 m; Noise = 3\%}
    	\end{center}
\end{frame}
\begin{frame}{Displaced source}
	\begin{center}
        	\includegraphics[height=0.75\textheight]{spostata_calibrata.pdf}
        	\vspace{0.3cm}
        	\scriptsize \textcolor{gray}{N true particles = 10; Pixel length = 0.00025 m; Noise = 3\%}
    	\end{center}    
\end{frame}
\begin{frame}{Displaced source}
	\begin{center}
        	\includegraphics[height=0.75\textheight]{spostata_distro.pdf}
        	\vspace{0.3cm}
        	\scriptsize \textcolor{gray}{Calibration with 100k events. N true particles = 10; Pixel length = 0.00025 m; Noise = 3\%}
    	\end{center}
\end{frame}
\begin{frame}{Displaced source}
	\begin{center}
        	\includegraphics[height=0.75\textheight]{spostata_distro_calibrata.pdf}
        	\vspace{0.3cm}
        	\scriptsize \textcolor{gray}{Calibration with 100k events. N true particles = 10; Pixel length = 0.00025 m; Noise = 3\%}
    	\end{center}
\end{frame}
\begin{frame}{Displaced source}
	\begin{center}
        	\includegraphics[height=0.75\textheight]{spostata_fitgauss.pdf}
        	\vspace{0.3cm}
        	\scriptsize \textcolor{gray}{Calibration with 100k events. N true particles = 10; Pixel length = 0.00025 m; Noise = 3\%}
    	\end{center}
\end{frame}
\begin{frame}{Displaced source}
	\begin{center}
        	\includegraphics[height=0.5\textheight]{spostata_accuratezza.pdf}
        	\vspace{0.3cm}
        	\scriptsize \textcolor{gray}{Calibration with 20k events per point. N true particles = 10; Pixel length = 0.00025 m; Noise = 3\%}
    	\end{center}
\end{frame}
\begin{frame}{Gaussian Generation}
	\begin{center}
        	\includegraphics[height=0.75\textheight]{gauss_distro.pdf}
        	\vspace{0.3cm}
        	\scriptsize \textcolor{gray}{Gaussian generation of the particle incidence angle. N true particles = 100; Pixel length = 0.0025 m; Noise = 0\%}
    	\end{center}
\end{frame}
\begin{frame}{Gaussian Generation}
	\begin{center}
        	\includegraphics[height=0.75\textheight]{gauss_distro_cali.pdf}
        	\vspace{0.3cm}
        	\scriptsize \textcolor{gray}{Gaussian generation of the particle incidence angle. N true particles = 100; Pixel length = 0.0025 m; Noise = 0\%}
    	\end{center}
\end{frame}




\end{document}
